\section{Conclusiones}

\begin{itemize}
  \item Con la cámara de un celular y un portátil se obtuvo un modelo 3D
        cerrado de un objeto en alrededor de un segundo, sin fondo especial y
        sin equipo adicional.
  \item El paso clave, ir de la imagen plana al espacio, se resolvió con la
        matriz de la cámara $K$ y su inversa: $\vect{P}=Z\,K^{-1}[u,v,1]^{T}$.
        Como la proyección deja un sistema de dos ecuaciones con tres
        incógnitas, la profundidad $Z$ tuvo que salir de información adicional:
        la silueta en el modo de una foto y las otras vistas en el modo de
        cinco fotos.
  \item Guardar la malla como dos matrices, $V$ ($N\times3$) y $F$
        ($M\times3$), permitió construirla, cerrarla, medir su volumen con
        determinantes y exportarla, todo con operaciones de matrices.
  \item Las matrices de rotación permitieron combinar cinco fotos tomadas
        desde lados distintos. Con las cinco vistas, las proporciones de una
        caja de prueba se recuperaron con menos del 3\,\% de error.
  \item Los temas básicos del curso (vectores, matrices y sistemas de
        ecuaciones) alcanzan para explicar y programar todo el proceso
        matemático de un problema real de visión por computador.
\end{itemize}
